% Figure captions as macros: % Figure captions as macros: % Figure captions as macros: % Figure captions as macros: \input{captions_figures} then use e.g. \CapFigSystemArchitecture
% Generated by paper_assets/scripts/make_docs.py
\newcommand{\CapFigZeroOneSystemArchitecture}{Overall architecture of the proposed system. The offline pipeline (top) trains a Random Forest on Backblaze SMART data (1 April -- 24 May 2019) and writes one prior-corrected 7-day failure score per drive-day for 1--23 June. Only these scores enter the trace-driven CloudSim simulation (bottom), where a risk classifier and an adaptive checkpoint scheduler set each host's checkpoint interval. Failure times come from a separate trace that only the failure injector reads; failure labels are never passed to the simulator. Orange borders: custom components; grey: native CloudSim.}
\newcommand{\CapFigZeroTwoMlPipeline}{Random Forest prediction pipeline. Steps 1--8 run offline; step 9 (risk classification) is applied inside the simulator. Labels use a 7-day horizon: a drive-day is positive if the drive records a failure on day t+1 to t+7. The tier boundaries are the 90th and 99th percentiles of validation scores and were not tuned on test data.}
\newcommand{\CapFigZeroThreeAdaptiveCheckpointing}{Adaptive checkpoint scheduling (policy E3). The Random Forest supplies only a risk score. A fixed classifier maps the score to a tier, and the scheduler maps the tier to a checkpoint interval $\tau$ relative to the baseline $T_0$. The table lists the resulting intervals for each $T_0$; the fixed policy E2 uses $\tau = T_0$ for every tier.}
\newcommand{\CapFigZeroFourCloudsimIntegration}{CloudSim integration and recovery architecture. Native CloudSim 3.0.3 entities (grey) are used unmodified. The custom broker runs each job as a chain of compute-segment cloudlets and writes each checkpoint as a 300 s cloudlet. On a host failure it moves the job to a spare host and restores from the last completed checkpoint. Each Backblaze drive is modelled as one host, and a recorded disk failure is treated as a host failure; this is a modelling assumption.}
\newcommand{\CapFigZeroFiveExecutionFlow}{End-to-end execution flow for one job. A job runs in compute segments that end at the next checkpoint due time, at midnight (when the daily risk tier may change) or at job end. A host failure interrupts the running segment or checkpoint write. The job then resumes on a spare host after $D + B = 180$ s: from its last completed checkpoint after a 300 s restore, or from zero if no checkpoint exists.}
\newcommand{\CapFigZeroSixExperimentalDesign}{Experimental design. In each scenario, all 11 policy configurations (E1, plus E2 and E3 at five baseline intervals $T_0$) run on all 56 shards with the identical workload, failure trace and risk trace, so comparisons are paired. The two load scenarios ($\rho = 0.9$ primary, $\rho = 1.0$ saturated sensitivity) have different workloads and are analysed separately.}
\newcommand{\CapGZeroOneRfRocCurve}{ROC curve of the Random Forest on the 1--23 June test set (2,544,754 drive-days, 997 positives). (a) Full curve, AUC = 0.804. (b) The same curve with a logarithmic FPR axis; markers show recall at 0.1\,\% FPR (16.0\,\%) and 1\,\% FPR (37.3\,\%).}
\newcommand{\CapGZeroTwoRfPrecisionRecallCurve}{Precision--recall curve on the June test set (PR-AUC = 0.0337, about 86 times the 0.039\,\% prevalence; log-scaled precision axis). The marker shows the alarm threshold 0.01368 chosen by maximum F1 on validation: precision 4.3\,\%, recall 20.7\,\%.}
\newcommand{\CapGZeroThreeRfConfusionMatrix}{Confusion matrix on the June test set at the validation-selected threshold ($p \ge 0.01368$). Shading shows the row-normalised rate. The model raises 4,597 false alarms for 206 true positives and misses 79\,\% of positive drive-days.}
\newcommand{\CapGZeroFourRfScoreDistribution}{Distribution of prior-corrected scores (log scale) for June drive-days with and without a failure in the next 7 days, each normalised to 100\,\%. Dotted lines mark the tier boundaries (0.000234, 0.001732) and the alarm threshold (0.01368). The two classes overlap substantially.}
\newcommand{\CapGZeroFiveRfFeatureImportance}{Top 15 features. (a) Mean decrease in impurity, final model. (b) Permutation importance (mean decrease in average precision $\pm$ s.d.) of the internal-validation model on a validation subsample. Reallocated-sector counts (SMART 5) dominate both rankings. Correlated features share importance, so the rankings should be read in groups.}
\newcommand{\CapGZeroSixOverheadVsTZeroRhoZeroNine}{Checkpoint overhead as a percentage of useful computation for fixed (E2) and RF-adaptive (E3) checkpointing at each baseline interval $T_0$, $\rho = 0.9$, all hosts (log scale; E1 has no overhead). Whiskers are 95\,\% host-clustered bootstrap CIs and are very narrow. E3 is lower at $T_0 \le 12$\,h and higher at 24--72\,h.}
\newcommand{\CapGZeroSevenLostWorkVsTZeroRhoZeroNine}{Lost computation caused by the 175 host failures, $\rho = 0.9$ (log scale). The dashed line shows no checkpointing (E1, 8,375\,h). Whiskers are 95\,\% bootstrap CIs; they are wide because they rest on 175 failure events.}
\newcommand{\CapGZeroEightDeadlineMissVsTZeroRhoZeroNine}{Deadline miss rate, $\rho = 0.9$. (a) All hosts: only E3 at $T_0 = 1$\,h misses an appreciable share (4.54\,\%). Every other E2/E3 configuration misses at most 0.015\,\%; value labels below 0.05\,\% are omitted. (b) The 175 failing hosts: E1 misses 19.8\,\%, and E3 at 1\,h misses 23.3\,\%. Whiskers are 95\,\% bootstrap CIs.}
\newcommand{\CapGZeroNineWastedTimeAllPoliciesRhoZeroNine}{Total wasted time (checkpoint overhead + lost work + interrupted checkpoint writes + recovery) as a percentage of useful computation, $\rho = 0.9$ (log scale). No checkpointing (E1) wastes the least overall (0.011\,\%). At these failure rates, checkpointing pays off by protecting the jobs that failures hit (Fig.~\ref{fig:deadline_rho09}b), not in aggregate efficiency.}
\newcommand{\CapGTenOverheadVsLostTradeoffRhoZeroNine}{Trade-off between overhead and lost computation, $\rho = 0.9$ (log--log). Each point is one policy at one $T_0$; labels are on the E2 points, and the E3 point for the same $T_0$ lies just below. E3 is lower on both axes at $T_0 \le 12$\,h; the 1\,h point carries a deadline penalty (Fig.~\ref{fig:deadline_rho09}). At 24 and 72\,h, E3 trades higher overhead for lower lost work. Whiskers: 95\,\% CIs of lost work.}
\newcommand{\CapGElevenCheckpointsPerJobRhoZeroNine}{Mean checkpoints per job, $\rho = 0.9$ (log scale). E3 writes fewer checkpoints than E2 at $T_0 \le 12$\,h, because 54\,\% of June host-days are LOW, and more at 24--72\,h.}
\newcommand{\CapGTwelvePairedEThreeVsETwoRhoZeroNine}{Paired E3 $-$ E2 differences at equal $T_0$, $\rho = 0.9$, all hosts: (a)--(c) relative to E2, (d) in percentage points. Intervals are 95\,\% host-clustered bootstrap CIs of the difference, rescaled by the E2 point estimate in (a)--(c). Filled markers: the interval excludes 0. Negative values favour E3. The intervals are not adjusted for multiple comparisons, and no significance is claimed.}
\newcommand{\CapGTwelvePairedEThreeVsETwoRhoOneZero}{Paired E3 $-$ E2 differences at equal $T_0$ under saturation, $\rho = 1.0$, all hosts. The format is the same as Fig.~\ref{fig:paired_rho09}. This experiment has a different workload from $\rho = 0.9$, so its effects are compared only in direction and size.}
\newcommand{\CapGThirteenDeadlineMissVsTZeroRhoOneZero}{Deadline miss rate under saturation, $\rho = 1.0$, all hosts. Without slack, checkpoint overhead becomes queueing delay: at $T_0 = 1$\,h, E2 misses 50.7\,\% of deadlines and E3 34.2\,\%; at 4\,h, 7.6\,\% versus 3.2\,\%. At 12\,h, E3 is worse (0.27\,\% versus 0.01\,\%).}
\newcommand{\CapGFourteenQueueingVsTZeroRhoOneZero}{Mean queueing delay per job, $\rho = 1.0$ (log scale). E3 reduces queueing at every $T_0$. The dashed line shows E1 (0.04\,h).}
\newcommand{\CapGFifteenLostWorkVsTZeroRhoOneZero}{Lost computation, $\rho = 1.0$ (log scale), with 95\,\% bootstrap CIs. The dashed line shows E1 (9,603\,h).}
\newcommand{\CapGSixteenFailingHostDeadlineMissBothRho}{Deadline miss rate on the 175 failing hosts, fixed (E2) versus adaptive (E3): (a) $\rho = 0.9$, (b) $\rho = 1.0$. Under saturation, E3's extra checkpoints on HIGH-tier hosts raise failing-host deadline misses at $T_0 \le 4$\,h (25.5\,\% versus 14.9\,\% at 4\,h), even though lost work falls. The two panels are separate experiments with different workloads.}
\newcommand{\CapGSeventeenTierOccupancy}{Share of June host-days in each risk tier, compared with the validation design target (90/9/1\,\%). June scores are shifted upward: only 54.1\,\% of host-days are LOW, and 6.6\,\% are HIGH. Failing hosts spend 52.3\,\% of their host-days in HIGH.}
\newcommand{\CapGEighteenFailuresByTier}{Risk tier in force at the instant of each of the 175 failures, identical across policies and scenarios: 110 HIGH, 38 MEDIUM and 27 LOW.}
\newcommand{\CapGNineteenLostWorkByTierRhoZeroNine}{Lost computation grouped by the tier in force at failure, $\rho = 0.9$ (log scale; point estimates). E3 sharply reduces lost work for failures in HIGH and has little effect in MEDIUM. It increases lost work for the 27 failures the model ranked LOW, where its interval is 2.65\,$T_0$.}
\newcommand{\CapGTwentyOverheadByTierRhoZeroNine}{Checkpoint overhead split by the tier in force when each checkpoint was written, $\rho = 0.9$, all hosts, indexed to the E2 total at the same $T_0$ (= 100). E3 cuts LOW-tier overhead and spends much of the saving on HIGH-tier checkpoints. At $T_0 = 72$\,h the LOW interval (190.8\,h) exceeds the longest job (168\,h), so E3 writes almost no LOW-tier checkpoints (0.2\,\% of its checkpoints).}
 then use e.g. \CapFigSystemArchitecture
% Generated by paper_assets/scripts/make_docs.py
\newcommand{\CapFigZeroOneSystemArchitecture}{Overall architecture of the proposed system. The offline pipeline (top) trains a Random Forest on Backblaze SMART data (1 April -- 24 May 2019) and writes one prior-corrected 7-day failure score per drive-day for 1--23 June. Only these scores enter the trace-driven CloudSim simulation (bottom), where a risk classifier and an adaptive checkpoint scheduler set each host's checkpoint interval. Failure times come from a separate trace that only the failure injector reads; failure labels are never passed to the simulator. Orange borders: custom components; grey: native CloudSim.}
\newcommand{\CapFigZeroTwoMlPipeline}{Random Forest prediction pipeline. Steps 1--8 run offline; step 9 (risk classification) is applied inside the simulator. Labels use a 7-day horizon: a drive-day is positive if the drive records a failure on day t+1 to t+7. The tier boundaries are the 90th and 99th percentiles of validation scores and were not tuned on test data.}
\newcommand{\CapFigZeroThreeAdaptiveCheckpointing}{Adaptive checkpoint scheduling (policy E3). The Random Forest supplies only a risk score. A fixed classifier maps the score to a tier, and the scheduler maps the tier to a checkpoint interval $\tau$ relative to the baseline $T_0$. The table lists the resulting intervals for each $T_0$; the fixed policy E2 uses $\tau = T_0$ for every tier.}
\newcommand{\CapFigZeroFourCloudsimIntegration}{CloudSim integration and recovery architecture. Native CloudSim 3.0.3 entities (grey) are used unmodified. The custom broker runs each job as a chain of compute-segment cloudlets and writes each checkpoint as a 300 s cloudlet. On a host failure it moves the job to a spare host and restores from the last completed checkpoint. Each Backblaze drive is modelled as one host, and a recorded disk failure is treated as a host failure; this is a modelling assumption.}
\newcommand{\CapFigZeroFiveExecutionFlow}{End-to-end execution flow for one job. A job runs in compute segments that end at the next checkpoint due time, at midnight (when the daily risk tier may change) or at job end. A host failure interrupts the running segment or checkpoint write. The job then resumes on a spare host after $D + B = 180$ s: from its last completed checkpoint after a 300 s restore, or from zero if no checkpoint exists.}
\newcommand{\CapFigZeroSixExperimentalDesign}{Experimental design. In each scenario, all 11 policy configurations (E1, plus E2 and E3 at five baseline intervals $T_0$) run on all 56 shards with the identical workload, failure trace and risk trace, so comparisons are paired. The two load scenarios ($\rho = 0.9$ primary, $\rho = 1.0$ saturated sensitivity) have different workloads and are analysed separately.}
\newcommand{\CapGZeroOneRfRocCurve}{ROC curve of the Random Forest on the 1--23 June test set (2,544,754 drive-days, 997 positives). (a) Full curve, AUC = 0.804. (b) The same curve with a logarithmic FPR axis; markers show recall at 0.1\,\% FPR (16.0\,\%) and 1\,\% FPR (37.3\,\%).}
\newcommand{\CapGZeroTwoRfPrecisionRecallCurve}{Precision--recall curve on the June test set (PR-AUC = 0.0337, about 86 times the 0.039\,\% prevalence; log-scaled precision axis). The marker shows the alarm threshold 0.01368 chosen by maximum F1 on validation: precision 4.3\,\%, recall 20.7\,\%.}
\newcommand{\CapGZeroThreeRfConfusionMatrix}{Confusion matrix on the June test set at the validation-selected threshold ($p \ge 0.01368$). Shading shows the row-normalised rate. The model raises 4,597 false alarms for 206 true positives and misses 79\,\% of positive drive-days.}
\newcommand{\CapGZeroFourRfScoreDistribution}{Distribution of prior-corrected scores (log scale) for June drive-days with and without a failure in the next 7 days, each normalised to 100\,\%. Dotted lines mark the tier boundaries (0.000234, 0.001732) and the alarm threshold (0.01368). The two classes overlap substantially.}
\newcommand{\CapGZeroFiveRfFeatureImportance}{Top 15 features. (a) Mean decrease in impurity, final model. (b) Permutation importance (mean decrease in average precision $\pm$ s.d.) of the internal-validation model on a validation subsample. Reallocated-sector counts (SMART 5) dominate both rankings. Correlated features share importance, so the rankings should be read in groups.}
\newcommand{\CapGZeroSixOverheadVsTZeroRhoZeroNine}{Checkpoint overhead as a percentage of useful computation for fixed (E2) and RF-adaptive (E3) checkpointing at each baseline interval $T_0$, $\rho = 0.9$, all hosts (log scale; E1 has no overhead). Whiskers are 95\,\% host-clustered bootstrap CIs and are very narrow. E3 is lower at $T_0 \le 12$\,h and higher at 24--72\,h.}
\newcommand{\CapGZeroSevenLostWorkVsTZeroRhoZeroNine}{Lost computation caused by the 175 host failures, $\rho = 0.9$ (log scale). The dashed line shows no checkpointing (E1, 8,375\,h). Whiskers are 95\,\% bootstrap CIs; they are wide because they rest on 175 failure events.}
\newcommand{\CapGZeroEightDeadlineMissVsTZeroRhoZeroNine}{Deadline miss rate, $\rho = 0.9$. (a) All hosts: only E3 at $T_0 = 1$\,h misses an appreciable share (4.54\,\%). Every other E2/E3 configuration misses at most 0.015\,\%; value labels below 0.05\,\% are omitted. (b) The 175 failing hosts: E1 misses 19.8\,\%, and E3 at 1\,h misses 23.3\,\%. Whiskers are 95\,\% bootstrap CIs.}
\newcommand{\CapGZeroNineWastedTimeAllPoliciesRhoZeroNine}{Total wasted time (checkpoint overhead + lost work + interrupted checkpoint writes + recovery) as a percentage of useful computation, $\rho = 0.9$ (log scale). No checkpointing (E1) wastes the least overall (0.011\,\%). At these failure rates, checkpointing pays off by protecting the jobs that failures hit (Fig.~\ref{fig:deadline_rho09}b), not in aggregate efficiency.}
\newcommand{\CapGTenOverheadVsLostTradeoffRhoZeroNine}{Trade-off between overhead and lost computation, $\rho = 0.9$ (log--log). Each point is one policy at one $T_0$; labels are on the E2 points, and the E3 point for the same $T_0$ lies just below. E3 is lower on both axes at $T_0 \le 12$\,h; the 1\,h point carries a deadline penalty (Fig.~\ref{fig:deadline_rho09}). At 24 and 72\,h, E3 trades higher overhead for lower lost work. Whiskers: 95\,\% CIs of lost work.}
\newcommand{\CapGElevenCheckpointsPerJobRhoZeroNine}{Mean checkpoints per job, $\rho = 0.9$ (log scale). E3 writes fewer checkpoints than E2 at $T_0 \le 12$\,h, because 54\,\% of June host-days are LOW, and more at 24--72\,h.}
\newcommand{\CapGTwelvePairedEThreeVsETwoRhoZeroNine}{Paired E3 $-$ E2 differences at equal $T_0$, $\rho = 0.9$, all hosts: (a)--(c) relative to E2, (d) in percentage points. Intervals are 95\,\% host-clustered bootstrap CIs of the difference, rescaled by the E2 point estimate in (a)--(c). Filled markers: the interval excludes 0. Negative values favour E3. The intervals are not adjusted for multiple comparisons, and no significance is claimed.}
\newcommand{\CapGTwelvePairedEThreeVsETwoRhoOneZero}{Paired E3 $-$ E2 differences at equal $T_0$ under saturation, $\rho = 1.0$, all hosts. The format is the same as Fig.~\ref{fig:paired_rho09}. This experiment has a different workload from $\rho = 0.9$, so its effects are compared only in direction and size.}
\newcommand{\CapGThirteenDeadlineMissVsTZeroRhoOneZero}{Deadline miss rate under saturation, $\rho = 1.0$, all hosts. Without slack, checkpoint overhead becomes queueing delay: at $T_0 = 1$\,h, E2 misses 50.7\,\% of deadlines and E3 34.2\,\%; at 4\,h, 7.6\,\% versus 3.2\,\%. At 12\,h, E3 is worse (0.27\,\% versus 0.01\,\%).}
\newcommand{\CapGFourteenQueueingVsTZeroRhoOneZero}{Mean queueing delay per job, $\rho = 1.0$ (log scale). E3 reduces queueing at every $T_0$. The dashed line shows E1 (0.04\,h).}
\newcommand{\CapGFifteenLostWorkVsTZeroRhoOneZero}{Lost computation, $\rho = 1.0$ (log scale), with 95\,\% bootstrap CIs. The dashed line shows E1 (9,603\,h).}
\newcommand{\CapGSixteenFailingHostDeadlineMissBothRho}{Deadline miss rate on the 175 failing hosts, fixed (E2) versus adaptive (E3): (a) $\rho = 0.9$, (b) $\rho = 1.0$. Under saturation, E3's extra checkpoints on HIGH-tier hosts raise failing-host deadline misses at $T_0 \le 4$\,h (25.5\,\% versus 14.9\,\% at 4\,h), even though lost work falls. The two panels are separate experiments with different workloads.}
\newcommand{\CapGSeventeenTierOccupancy}{Share of June host-days in each risk tier, compared with the validation design target (90/9/1\,\%). June scores are shifted upward: only 54.1\,\% of host-days are LOW, and 6.6\,\% are HIGH. Failing hosts spend 52.3\,\% of their host-days in HIGH.}
\newcommand{\CapGEighteenFailuresByTier}{Risk tier in force at the instant of each of the 175 failures, identical across policies and scenarios: 110 HIGH, 38 MEDIUM and 27 LOW.}
\newcommand{\CapGNineteenLostWorkByTierRhoZeroNine}{Lost computation grouped by the tier in force at failure, $\rho = 0.9$ (log scale; point estimates). E3 sharply reduces lost work for failures in HIGH and has little effect in MEDIUM. It increases lost work for the 27 failures the model ranked LOW, where its interval is 2.65\,$T_0$.}
\newcommand{\CapGTwentyOverheadByTierRhoZeroNine}{Checkpoint overhead split by the tier in force when each checkpoint was written, $\rho = 0.9$, all hosts, indexed to the E2 total at the same $T_0$ (= 100). E3 cuts LOW-tier overhead and spends much of the saving on HIGH-tier checkpoints. At $T_0 = 72$\,h the LOW interval (190.8\,h) exceeds the longest job (168\,h), so E3 writes almost no LOW-tier checkpoints (0.2\,\% of its checkpoints).}
 then use e.g. \CapFigSystemArchitecture
% Generated by paper_assets/scripts/make_docs.py
\newcommand{\CapFigZeroOneSystemArchitecture}{Overall architecture of the proposed system. The offline pipeline (top) trains a Random Forest on Backblaze SMART data (1 April -- 24 May 2019) and writes one prior-corrected 7-day failure score per drive-day for 1--23 June. Only these scores enter the trace-driven CloudSim simulation (bottom), where a risk classifier and an adaptive checkpoint scheduler set each host's checkpoint interval. Failure times come from a separate trace that only the failure injector reads; failure labels are never passed to the simulator. Orange borders: custom components; grey: native CloudSim.}
\newcommand{\CapFigZeroTwoMlPipeline}{Random Forest prediction pipeline. Steps 1--8 run offline; step 9 (risk classification) is applied inside the simulator. Labels use a 7-day horizon: a drive-day is positive if the drive records a failure on day t+1 to t+7. The tier boundaries are the 90th and 99th percentiles of validation scores and were not tuned on test data.}
\newcommand{\CapFigZeroThreeAdaptiveCheckpointing}{Adaptive checkpoint scheduling (policy E3). The Random Forest supplies only a risk score. A fixed classifier maps the score to a tier, and the scheduler maps the tier to a checkpoint interval $\tau$ relative to the baseline $T_0$. The table lists the resulting intervals for each $T_0$; the fixed policy E2 uses $\tau = T_0$ for every tier.}
\newcommand{\CapFigZeroFourCloudsimIntegration}{CloudSim integration and recovery architecture. Native CloudSim 3.0.3 entities (grey) are used unmodified. The custom broker runs each job as a chain of compute-segment cloudlets and writes each checkpoint as a 300 s cloudlet. On a host failure it moves the job to a spare host and restores from the last completed checkpoint. Each Backblaze drive is modelled as one host, and a recorded disk failure is treated as a host failure; this is a modelling assumption.}
\newcommand{\CapFigZeroFiveExecutionFlow}{End-to-end execution flow for one job. A job runs in compute segments that end at the next checkpoint due time, at midnight (when the daily risk tier may change) or at job end. A host failure interrupts the running segment or checkpoint write. The job then resumes on a spare host after $D + B = 180$ s: from its last completed checkpoint after a 300 s restore, or from zero if no checkpoint exists.}
\newcommand{\CapFigZeroSixExperimentalDesign}{Experimental design. In each scenario, all 11 policy configurations (E1, plus E2 and E3 at five baseline intervals $T_0$) run on all 56 shards with the identical workload, failure trace and risk trace, so comparisons are paired. The two load scenarios ($\rho = 0.9$ primary, $\rho = 1.0$ saturated sensitivity) have different workloads and are analysed separately.}
\newcommand{\CapGZeroOneRfRocCurve}{ROC curve of the Random Forest on the 1--23 June test set (2,544,754 drive-days, 997 positives). (a) Full curve, AUC = 0.804. (b) The same curve with a logarithmic FPR axis; markers show recall at 0.1\,\% FPR (16.0\,\%) and 1\,\% FPR (37.3\,\%).}
\newcommand{\CapGZeroTwoRfPrecisionRecallCurve}{Precision--recall curve on the June test set (PR-AUC = 0.0337, about 86 times the 0.039\,\% prevalence; log-scaled precision axis). The marker shows the alarm threshold 0.01368 chosen by maximum F1 on validation: precision 4.3\,\%, recall 20.7\,\%.}
\newcommand{\CapGZeroThreeRfConfusionMatrix}{Confusion matrix on the June test set at the validation-selected threshold ($p \ge 0.01368$). Shading shows the row-normalised rate. The model raises 4,597 false alarms for 206 true positives and misses 79\,\% of positive drive-days.}
\newcommand{\CapGZeroFourRfScoreDistribution}{Distribution of prior-corrected scores (log scale) for June drive-days with and without a failure in the next 7 days, each normalised to 100\,\%. Dotted lines mark the tier boundaries (0.000234, 0.001732) and the alarm threshold (0.01368). The two classes overlap substantially.}
\newcommand{\CapGZeroFiveRfFeatureImportance}{Top 15 features. (a) Mean decrease in impurity, final model. (b) Permutation importance (mean decrease in average precision $\pm$ s.d.) of the internal-validation model on a validation subsample. Reallocated-sector counts (SMART 5) dominate both rankings. Correlated features share importance, so the rankings should be read in groups.}
\newcommand{\CapGZeroSixOverheadVsTZeroRhoZeroNine}{Checkpoint overhead as a percentage of useful computation for fixed (E2) and RF-adaptive (E3) checkpointing at each baseline interval $T_0$, $\rho = 0.9$, all hosts (log scale; E1 has no overhead). Whiskers are 95\,\% host-clustered bootstrap CIs and are very narrow. E3 is lower at $T_0 \le 12$\,h and higher at 24--72\,h.}
\newcommand{\CapGZeroSevenLostWorkVsTZeroRhoZeroNine}{Lost computation caused by the 175 host failures, $\rho = 0.9$ (log scale). The dashed line shows no checkpointing (E1, 8,375\,h). Whiskers are 95\,\% bootstrap CIs; they are wide because they rest on 175 failure events.}
\newcommand{\CapGZeroEightDeadlineMissVsTZeroRhoZeroNine}{Deadline miss rate, $\rho = 0.9$. (a) All hosts: only E3 at $T_0 = 1$\,h misses an appreciable share (4.54\,\%). Every other E2/E3 configuration misses at most 0.015\,\%; value labels below 0.05\,\% are omitted. (b) The 175 failing hosts: E1 misses 19.8\,\%, and E3 at 1\,h misses 23.3\,\%. Whiskers are 95\,\% bootstrap CIs.}
\newcommand{\CapGZeroNineWastedTimeAllPoliciesRhoZeroNine}{Total wasted time (checkpoint overhead + lost work + interrupted checkpoint writes + recovery) as a percentage of useful computation, $\rho = 0.9$ (log scale). No checkpointing (E1) wastes the least overall (0.011\,\%). At these failure rates, checkpointing pays off by protecting the jobs that failures hit (Fig.~\ref{fig:deadline_rho09}b), not in aggregate efficiency.}
\newcommand{\CapGTenOverheadVsLostTradeoffRhoZeroNine}{Trade-off between overhead and lost computation, $\rho = 0.9$ (log--log). Each point is one policy at one $T_0$; labels are on the E2 points, and the E3 point for the same $T_0$ lies just below. E3 is lower on both axes at $T_0 \le 12$\,h; the 1\,h point carries a deadline penalty (Fig.~\ref{fig:deadline_rho09}). At 24 and 72\,h, E3 trades higher overhead for lower lost work. Whiskers: 95\,\% CIs of lost work.}
\newcommand{\CapGElevenCheckpointsPerJobRhoZeroNine}{Mean checkpoints per job, $\rho = 0.9$ (log scale). E3 writes fewer checkpoints than E2 at $T_0 \le 12$\,h, because 54\,\% of June host-days are LOW, and more at 24--72\,h.}
\newcommand{\CapGTwelvePairedEThreeVsETwoRhoZeroNine}{Paired E3 $-$ E2 differences at equal $T_0$, $\rho = 0.9$, all hosts: (a)--(c) relative to E2, (d) in percentage points. Intervals are 95\,\% host-clustered bootstrap CIs of the difference, rescaled by the E2 point estimate in (a)--(c). Filled markers: the interval excludes 0. Negative values favour E3. The intervals are not adjusted for multiple comparisons, and no significance is claimed.}
\newcommand{\CapGTwelvePairedEThreeVsETwoRhoOneZero}{Paired E3 $-$ E2 differences at equal $T_0$ under saturation, $\rho = 1.0$, all hosts. The format is the same as Fig.~\ref{fig:paired_rho09}. This experiment has a different workload from $\rho = 0.9$, so its effects are compared only in direction and size.}
\newcommand{\CapGThirteenDeadlineMissVsTZeroRhoOneZero}{Deadline miss rate under saturation, $\rho = 1.0$, all hosts. Without slack, checkpoint overhead becomes queueing delay: at $T_0 = 1$\,h, E2 misses 50.7\,\% of deadlines and E3 34.2\,\%; at 4\,h, 7.6\,\% versus 3.2\,\%. At 12\,h, E3 is worse (0.27\,\% versus 0.01\,\%).}
\newcommand{\CapGFourteenQueueingVsTZeroRhoOneZero}{Mean queueing delay per job, $\rho = 1.0$ (log scale). E3 reduces queueing at every $T_0$. The dashed line shows E1 (0.04\,h).}
\newcommand{\CapGFifteenLostWorkVsTZeroRhoOneZero}{Lost computation, $\rho = 1.0$ (log scale), with 95\,\% bootstrap CIs. The dashed line shows E1 (9,603\,h).}
\newcommand{\CapGSixteenFailingHostDeadlineMissBothRho}{Deadline miss rate on the 175 failing hosts, fixed (E2) versus adaptive (E3): (a) $\rho = 0.9$, (b) $\rho = 1.0$. Under saturation, E3's extra checkpoints on HIGH-tier hosts raise failing-host deadline misses at $T_0 \le 4$\,h (25.5\,\% versus 14.9\,\% at 4\,h), even though lost work falls. The two panels are separate experiments with different workloads.}
\newcommand{\CapGSeventeenTierOccupancy}{Share of June host-days in each risk tier, compared with the validation design target (90/9/1\,\%). June scores are shifted upward: only 54.1\,\% of host-days are LOW, and 6.6\,\% are HIGH. Failing hosts spend 52.3\,\% of their host-days in HIGH.}
\newcommand{\CapGEighteenFailuresByTier}{Risk tier in force at the instant of each of the 175 failures, identical across policies and scenarios: 110 HIGH, 38 MEDIUM and 27 LOW.}
\newcommand{\CapGNineteenLostWorkByTierRhoZeroNine}{Lost computation grouped by the tier in force at failure, $\rho = 0.9$ (log scale; point estimates). E3 sharply reduces lost work for failures in HIGH and has little effect in MEDIUM. It increases lost work for the 27 failures the model ranked LOW, where its interval is 2.65\,$T_0$.}
\newcommand{\CapGTwentyOverheadByTierRhoZeroNine}{Checkpoint overhead split by the tier in force when each checkpoint was written, $\rho = 0.9$, all hosts, indexed to the E2 total at the same $T_0$ (= 100). E3 cuts LOW-tier overhead and spends much of the saving on HIGH-tier checkpoints. At $T_0 = 72$\,h the LOW interval (190.8\,h) exceeds the longest job (168\,h), so E3 writes almost no LOW-tier checkpoints (0.2\,\% of its checkpoints).}
 then use e.g. \CapFigSystemArchitecture
% Generated by paper_assets/scripts/make_docs.py
\newcommand{\CapFigZeroOneSystemArchitecture}{Overall architecture of the proposed system. The offline pipeline (top) trains a Random Forest on Backblaze SMART data (1 April -- 24 May 2019) and writes one prior-corrected 7-day failure score per drive-day for 1--23 June. Only these scores enter the trace-driven CloudSim simulation (bottom), where a risk classifier and an adaptive checkpoint scheduler set each host's checkpoint interval. Failure times come from a separate trace that only the failure injector reads; failure labels are never passed to the simulator. Orange borders: custom components; grey: native CloudSim.}
\newcommand{\CapFigZeroTwoMlPipeline}{Random Forest prediction pipeline. Steps 1--8 run offline; step 9 (risk classification) is applied inside the simulator. Labels use a 7-day horizon: a drive-day is positive if the drive records a failure on day t+1 to t+7. The tier boundaries are the 90th and 99th percentiles of validation scores and were not tuned on test data.}
\newcommand{\CapFigZeroThreeAdaptiveCheckpointing}{Adaptive checkpoint scheduling (policy E3). The Random Forest supplies only a risk score. A fixed classifier maps the score to a tier, and the scheduler maps the tier to a checkpoint interval $\tau$ relative to the baseline $T_0$. The table lists the resulting intervals for each $T_0$; the fixed policy E2 uses $\tau = T_0$ for every tier.}
\newcommand{\CapFigZeroFourCloudsimIntegration}{CloudSim integration and recovery architecture. Native CloudSim 3.0.3 entities (grey) are used unmodified. The custom broker runs each job as a chain of compute-segment cloudlets and writes each checkpoint as a 300 s cloudlet. On a host failure it moves the job to a spare host and restores from the last completed checkpoint. Each Backblaze drive is modelled as one host, and a recorded disk failure is treated as a host failure; this is a modelling assumption.}
\newcommand{\CapFigZeroFiveExecutionFlow}{End-to-end execution flow for one job. A job runs in compute segments that end at the next checkpoint due time, at midnight (when the daily risk tier may change) or at job end. A host failure interrupts the running segment or checkpoint write. The job then resumes on a spare host after $D + B = 180$ s: from its last completed checkpoint after a 300 s restore, or from zero if no checkpoint exists.}
\newcommand{\CapFigZeroSixExperimentalDesign}{Experimental design. In each scenario, all 11 policy configurations (E1, plus E2 and E3 at five baseline intervals $T_0$) run on all 56 shards with the identical workload, failure trace and risk trace, so comparisons are paired. The two load scenarios ($\rho = 0.9$ primary, $\rho = 1.0$ saturated sensitivity) have different workloads and are analysed separately.}
\newcommand{\CapGZeroOneRfRocCurve}{ROC curve of the Random Forest on the 1--23 June test set (2,544,754 drive-days, 997 positives). (a) Full curve, AUC = 0.804. (b) The same curve with a logarithmic FPR axis; markers show recall at 0.1\,\% FPR (16.0\,\%) and 1\,\% FPR (37.3\,\%).}
\newcommand{\CapGZeroTwoRfPrecisionRecallCurve}{Precision--recall curve on the June test set (PR-AUC = 0.0337, about 86 times the 0.039\,\% prevalence; log-scaled precision axis). The marker shows the alarm threshold 0.01368 chosen by maximum F1 on validation: precision 4.3\,\%, recall 20.7\,\%.}
\newcommand{\CapGZeroThreeRfConfusionMatrix}{Confusion matrix on the June test set at the validation-selected threshold ($p \ge 0.01368$). Shading shows the row-normalised rate. The model raises 4,597 false alarms for 206 true positives and misses 79\,\% of positive drive-days.}
\newcommand{\CapGZeroFourRfScoreDistribution}{Distribution of prior-corrected scores (log scale) for June drive-days with and without a failure in the next 7 days, each normalised to 100\,\%. Dotted lines mark the tier boundaries (0.000234, 0.001732) and the alarm threshold (0.01368). The two classes overlap substantially.}
\newcommand{\CapGZeroFiveRfFeatureImportance}{Top 15 features. (a) Mean decrease in impurity, final model. (b) Permutation importance (mean decrease in average precision $\pm$ s.d.) of the internal-validation model on a validation subsample. Reallocated-sector counts (SMART 5) dominate both rankings. Correlated features share importance, so the rankings should be read in groups.}
\newcommand{\CapGZeroSixOverheadVsTZeroRhoZeroNine}{Checkpoint overhead as a percentage of useful computation for fixed (E2) and RF-adaptive (E3) checkpointing at each baseline interval $T_0$, $\rho = 0.9$, all hosts (log scale; E1 has no overhead). Whiskers are 95\,\% host-clustered bootstrap CIs and are very narrow. E3 is lower at $T_0 \le 12$\,h and higher at 24--72\,h.}
\newcommand{\CapGZeroSevenLostWorkVsTZeroRhoZeroNine}{Lost computation caused by the 175 host failures, $\rho = 0.9$ (log scale). The dashed line shows no checkpointing (E1, 8,375\,h). Whiskers are 95\,\% bootstrap CIs; they are wide because they rest on 175 failure events.}
\newcommand{\CapGZeroEightDeadlineMissVsTZeroRhoZeroNine}{Deadline miss rate, $\rho = 0.9$. (a) All hosts: only E3 at $T_0 = 1$\,h misses an appreciable share (4.54\,\%). Every other E2/E3 configuration misses at most 0.015\,\%; value labels below 0.05\,\% are omitted. (b) The 175 failing hosts: E1 misses 19.8\,\%, and E3 at 1\,h misses 23.3\,\%. Whiskers are 95\,\% bootstrap CIs.}
\newcommand{\CapGZeroNineWastedTimeAllPoliciesRhoZeroNine}{Total wasted time (checkpoint overhead + lost work + interrupted checkpoint writes + recovery) as a percentage of useful computation, $\rho = 0.9$ (log scale). No checkpointing (E1) wastes the least overall (0.011\,\%). At these failure rates, checkpointing pays off by protecting the jobs that failures hit (Fig.~\ref{fig:deadline_rho09}b), not in aggregate efficiency.}
\newcommand{\CapGTenOverheadVsLostTradeoffRhoZeroNine}{Trade-off between overhead and lost computation, $\rho = 0.9$ (log--log). Each point is one policy at one $T_0$; labels are on the E2 points, and the E3 point for the same $T_0$ lies just below. E3 is lower on both axes at $T_0 \le 12$\,h; the 1\,h point carries a deadline penalty (Fig.~\ref{fig:deadline_rho09}). At 24 and 72\,h, E3 trades higher overhead for lower lost work. Whiskers: 95\,\% CIs of lost work.}
\newcommand{\CapGElevenCheckpointsPerJobRhoZeroNine}{Mean checkpoints per job, $\rho = 0.9$ (log scale). E3 writes fewer checkpoints than E2 at $T_0 \le 12$\,h, because 54\,\% of June host-days are LOW, and more at 24--72\,h.}
\newcommand{\CapGTwelvePairedEThreeVsETwoRhoZeroNine}{Paired E3 $-$ E2 differences at equal $T_0$, $\rho = 0.9$, all hosts: (a)--(c) relative to E2, (d) in percentage points. Intervals are 95\,\% host-clustered bootstrap CIs of the difference, rescaled by the E2 point estimate in (a)--(c). Filled markers: the interval excludes 0. Negative values favour E3. The intervals are not adjusted for multiple comparisons, and no significance is claimed.}
\newcommand{\CapGTwelvePairedEThreeVsETwoRhoOneZero}{Paired E3 $-$ E2 differences at equal $T_0$ under saturation, $\rho = 1.0$, all hosts. The format is the same as Fig.~\ref{fig:paired_rho09}. This experiment has a different workload from $\rho = 0.9$, so its effects are compared only in direction and size.}
\newcommand{\CapGThirteenDeadlineMissVsTZeroRhoOneZero}{Deadline miss rate under saturation, $\rho = 1.0$, all hosts. Without slack, checkpoint overhead becomes queueing delay: at $T_0 = 1$\,h, E2 misses 50.7\,\% of deadlines and E3 34.2\,\%; at 4\,h, 7.6\,\% versus 3.2\,\%. At 12\,h, E3 is worse (0.27\,\% versus 0.01\,\%).}
\newcommand{\CapGFourteenQueueingVsTZeroRhoOneZero}{Mean queueing delay per job, $\rho = 1.0$ (log scale). E3 reduces queueing at every $T_0$. The dashed line shows E1 (0.04\,h).}
\newcommand{\CapGFifteenLostWorkVsTZeroRhoOneZero}{Lost computation, $\rho = 1.0$ (log scale), with 95\,\% bootstrap CIs. The dashed line shows E1 (9,603\,h).}
\newcommand{\CapGSixteenFailingHostDeadlineMissBothRho}{Deadline miss rate on the 175 failing hosts, fixed (E2) versus adaptive (E3): (a) $\rho = 0.9$, (b) $\rho = 1.0$. Under saturation, E3's extra checkpoints on HIGH-tier hosts raise failing-host deadline misses at $T_0 \le 4$\,h (25.5\,\% versus 14.9\,\% at 4\,h), even though lost work falls. The two panels are separate experiments with different workloads.}
\newcommand{\CapGSeventeenTierOccupancy}{Share of June host-days in each risk tier, compared with the validation design target (90/9/1\,\%). June scores are shifted upward: only 54.1\,\% of host-days are LOW, and 6.6\,\% are HIGH. Failing hosts spend 52.3\,\% of their host-days in HIGH.}
\newcommand{\CapGEighteenFailuresByTier}{Risk tier in force at the instant of each of the 175 failures, identical across policies and scenarios: 110 HIGH, 38 MEDIUM and 27 LOW.}
\newcommand{\CapGNineteenLostWorkByTierRhoZeroNine}{Lost computation grouped by the tier in force at failure, $\rho = 0.9$ (log scale; point estimates). E3 sharply reduces lost work for failures in HIGH and has little effect in MEDIUM. It increases lost work for the 27 failures the model ranked LOW, where its interval is 2.65\,$T_0$.}
\newcommand{\CapGTwentyOverheadByTierRhoZeroNine}{Checkpoint overhead split by the tier in force when each checkpoint was written, $\rho = 0.9$, all hosts, indexed to the E2 total at the same $T_0$ (= 100). E3 cuts LOW-tier overhead and spends much of the saving on HIGH-tier checkpoints. At $T_0 = 72$\,h the LOW interval (190.8\,h) exceeds the longest job (168\,h), so E3 writes almost no LOW-tier checkpoints (0.2\,\% of its checkpoints).}
